% ============================================================
%  DAFTAR PUSTAKA
%  Pedoman: alfabetik, spasi 1 per entri, 2 spasi antar entri
%  Format FTI Gunadarma
% ============================================================

\phantomsection
\chapter*{DAFTAR PUSTAKA}
\addcontentsline{toc}{chapter}{DAFTAR PUSTAKA}
\thispagestyle{chapterpage}

\printbibliography[heading=none]

%% ---- Buku ----
%Munawar. (2018). \textbf{Analisis Perancangan Sistem Berorientasi Objek dengan UML}.
%ISBN 9786026232779. Jakarta: Penerbit Informatika.
%
%% ---- Buku Online ----
%Budiman, Edy. (2018). \textbf{Mobile Programming for Student}. Mulawarman University Press.
%Diakses dari \url{https://www.academia.edu/38253550/Buku_Ajar_Mobile_Programing_for_Student.pdf}.
%
%% ---- Jurnal ----
%Budiman, Edy. (2016). Pemanfaatan Teknologi \textit{Location Based Service} dalam
%Pengembangan Aplikasi Profil Kampus Universitas Mulawarman Berbasis \textit{Mobile}.
%\textbf{Jurnal Ilmiah ILKOM} Volume 8 Nomor 3 (Desember 2016). ISSN: 2087-1716.
%
%% ---- Jurnal Online ----
%Barros, Bonifacio., Marisa, Fitri., dan Wijaya, Indra Dharma. (2018). Pembuatan
%\textit{Game} Kuis Siapa Pintar. \textbf{Jurnal Informatika Merdeka Pasuruan}.
%Vol 3, No 1 2018. Diakses tanggal 14 Februari 2020 dari
%\url{https://ejurnal.unmerpas.ac.id/index.php/informatika/article/view/88}.
%
%% ---- Artikel Online ----
%Nama Penulis. (Tanggal). Judul Artikel. Diakses dari \url{https://contoh.url/artikel}.
%
%% ---- Publikasi Pemerintah Online ----
%Kemenristekdikti. (2016). Sosialisasi Kemampuan Lulusan. Diakses dari
%\url{http://www.kopertis12.or.id/wp-content/uploads/2016/06/small_SN_DIKTI_44_2015_SOSIALISASI.pdf}.


%    \makeatletter
%    \renewcommand\@biblabel[1]{%
%        \textcolor{white}{[#1]}%  % sama dengan warna background
%    }
%    \makeatother
%
%    \begin{thebibliography}{99}
%
%        \bibitem{chung2020visual}
%        Chung, S. dan Son, J.-W. (2020). \textbf{Visual Perception in Autism Spectrum Disorder: A Review of Neuroimaging Studies}.
%        \textbf{Soa Chongsonyon Chongsin Uihak}, Volume 31 (2020), hlm. 105--120.
%
%        \bibitem{patterson2022color}
%        Patterson, E., Mastey, R., Kuchenbecker, J., Rowlan, J., Neitz, J., Neitz, M., dan Carroll, J. (2022). \textbf{Effects of color-enhancing glasses on color vision in congenital red-green color deficiencies}.
%        \textbf{Optics Express}, (2022), hlm. 31182--31194.
%
%        \bibitem{improving_discrimination}
%        Lin, Huei-Yung, Chen, Li-Qi, dan Wang, Min-Liang. (2019). \textbf{Improving Discrimination in Color Vision Deficiency by Image Re-Coloring}.
%        \textbf{Sensors}, Volume 19, Nomor 10, Article 2250. ISSN: 1424-8220.
%        Diakses dari \url{https://www.mdpi.com/1424-8220/19/10/2250}.
%
%        \bibitem{Elrefaei2018}
%        Elrefaei, Lamiaa A. (2018). \textbf{Smartphone Based Image Color Correction for Color Blindness}.
%        \textbf{International Journal of Interactive Mobile Technologies (iJIM)}, Volume 12, Nomor 3 (Juli 2018), hlm. 104--119.
%        Diakses dari \url{https://online-journals.org/index.php/i-jim/article/view/8160}.
%
%        \bibitem{Vinot1999DigitalVC}
%        Vi{\'e}not, Fran\c{c}oise, Brettel, Hans, dan Mollon, John D. (1999). \textbf{Digital video colourmaps for checking the legibility of displays by dichromats}.
%        \textbf{Color Research and Application}, Volume 24 (1999), hlm. 243--252.
%        Diakses dari \url{https://api.semanticscholar.org/CorpusID:62166956}.
%
%        \bibitem{muhammad2020analysis}
%        Muhammad, Aisha Badamasi, Maitama, Ja'afar Zubairu, Bature, Amir, dan Yunusa, Zainab. (2020). \textbf{Analysis of {LMS} Daltonization Algorithm for Adjusting Image Color for Deuteranopia}.
%        \textbf{Zaria Journal of Electrical Engineering Technology}, Volume 9, Nomor 1 (Maret 2020). ISSN: 0261-1570, hlm. 128--135.
%
%        \bibitem{tasnim2017}
%        Tasnim, Anika dan Hasan, Mohammad S. (2017). \textbf{An improved dynamic daltonization for color-blinds}.
%        Dalam \textbf{2017 IEEE Region 10 Humanitarian Technology Conference (R10-HTC)}, hlm. 798--801.
%
%        \bibitem{physiologically_visualization}
%        Machado, Gustavo M., Oliveira, Manuel M., dan Fernandes, Leandro A. F. (2009). \textbf{A Physiologically-based Model for Simulation of Color Vision Deficiency}.
%        \textbf{IEEE Transactions on Visualization and Computer Graphics}, Volume 15, Nomor 6 (2009), hlm. 1291--1298.
%
%        \bibitem{schmitz_colorblindness}
%        Schmitz, James. (2026). \textbf{Color Blindness Simulation Research}.
%        Diakses tanggal 26 Maret 2026 dari \url{https://ixora.io/projects/colorblindness/color-blindness-simulation-research/}.
%
%        \bibitem{statcounter2025mobile}
%        StatCounter. (2025). \textbf{Mobile Operating System Market Share Worldwide}.
%        Diakses tanggal 26 Februari 2026 dari \url{https://gs.statcounter.com/os-market-share/mobile/worldwide/2025}.
%
%        \bibitem{covisance}
%        Tecson, Gerald, Calanda, Fredilyn, Cayabyab, Gerald, dan Reyes Jr, Felizardo. (2017). \textbf{Covisance: A Real Time Mobile Recolorization Tool for Aiding Color Vision Deficient Users Utilizing D-15 Color Arrangement Test}.
%        Dalam \textbf{Proceedings of the 2017 ACM}, hlm. 95--99.
%
%        \bibitem{camerax_architecture}
%        Android Developers. (2024). \textbf{CameraX Architecture}.
%        Diakses tanggal 23 Maret 2026 dari \url{https://developer.android.com/media/camera/camerax/architecture}.
%
%        \bibitem{bci_biomedical}
%        Man, Dariusz dan Olchawa, Ryszard. (2018). \textbf{The Possibilities of Using BCI Technology in Biomedical Engineering}.
%        Dalam \textbf{Advances in Intelligent Systems and Computing}, Volume 720, hlm. 30--37.
%
%    \end{thebibliography}
%
%} % end singlespacing
